\section{해석과 한계}

Formal state는 서로 결합된 다섯 structural observable과 하나의 temporal
transition layer를 갖는다. Topology는 adjacency와 connectedness 같은 qualitative
organization을, geometry는 metric 정보를 다룬다. Relation은 typed entity,
generator morphism과 composition을 다룬다. Filtration과 mass는 같은 carrier에
ordered observable과 measure-theoretic observable을 추가한다. Dynamics는
action-conditioned tracked transition, history, turnover, intervention과 rollout을
다룬다. Common tagged carrier와 compatibility axiom은 이 layer들을 동일시하지
않으면서 결합한다.

이 형식 결과는 명시적으로 지정된 object에 관한 conditional statement다. 이
observable들이 intelligence에 필요하거나 충분하다는 것, human cognition이 이
state object를 그대로 사용한다는 것, unconstrained neural network가 이를
발견한다는 것을 증명하지 않는다. Structurally specified model이 보편적으로
우월하다는 것도 증명하지 않는다. 이러한 주장은 별도의 cognitive·empirical
validation을 요구한다.

핵심 design boundary는 의도적이다. Structural claim은 평가되기 전에 carrier,
equivalence, observable, transformation, discrepancy를 명시해야 한다. Scalar predictive
loss는 유용할 수 있지만 두 state가 structurally equivalent한지, transition이 필요한
relation을 보존하는지는 그것만으로 결정할 수 없다. 반대로 structural discrepancy도
task-relevant alignment와 invariant가 지정된 뒤에만 의미를 갖는다.

이 formalism은 entity set, 그 typing과 시간에 걸친 identity를 state의 주어진 data로
취급한다. Raw observation에서 entity를 발견하거나 perception error 아래에서 identity를
유지하는 procedure는 제공하지 않는다. 그러한 단계는 여기서 명시한 모든 것에 선행한다.

동반 empirical paper들은 이 framework의 제한된 consequence를 finite mechanism
family에서 검정한다. 이 실험들은 전체 이론이나 real-world cognition의
validation이 아니라 명시된 consequence의 validation으로 해석해야 한다.
